% !TeX root = thesis.tex
\chapter*{Glossary}
\phantomsection
\addcontentsline{toc}{chapter}{Glossary}

\noindent This section gives short explanations of important terms used throughout the thesis.

\begin{description}
\item[Biological graph] A directed graph built from prior molecular interactions, used to define how information can move between genes in the proposed model.
\item[Cell context] A sample specific vector derived from the matched control profile. It summarizes the baseline state of the cell before drug treatment.
\item[Cell identity embedding] A learned vector assigned to each cell line. All samples from the same cell line share the same embedding.
\item[Warm split] Evaluation split in which training and test samples come from similar observed conditions, so no major entity type is fully held out.
\item[Cold cell split] Evaluation split in which the test cell lines are not seen during training.
\item[Cold drug target split] Evaluation split in which the test drugs are held out from training to test generalization to unseen drug target patterns.
\item[Counterfactual drug loss] Training term that compares the normal drug input with a zero drug version in order to encourage the model to rely on the perturbation signal.
\item[Drug target vector] Binary vector that marks which genes are annotated targets of the drug.
\item[Landmark genes] The measured subset of 978 genes used in the LINCS L1000 assay.
\item[Perturbation] A change applied to a biological system, here caused by drug treatment.
\item[Response profile] The vector of gene expression changes caused by the drug, relative to the matched control profile.
\item[Uncertainty output] The predicted variance or log variance that estimates how confident the model is about each gene level prediction.
\end{description}
